\section{Background}
\label{sec:background}

This section sets up the vocabulary and the prior art on which our work
builds. We focus on the three concepts that structure the rest of the paper:
sampling strategies for configurable systems
(\Cref{subsec:sampling-strategies}), the cost models that have been associated
with them (\Cref{subsec:cost-models}), and what is and is not known about
context sensitivity of sampling-based conclusions
(\Cref{subsec:context-sensitivity}). \Cref{sec:related_work} returns to the
literature once we have introduced our own formulation.

\subsection{Configurable Systems and Performance Prediction}
\label{subsec:configurable-systems}

A configurable system exposes a set of options whose joint configuration affects
non-functional properties such as runtime, memory consumption, or output
size. Let $C$ be the set of \emph{valid} configurations, \ie those satisfying
constraints among options. For a fixed input $i$ and property $p$, every
$c\in C$ has a measured value $y_{c}\in\mathbb{R}$ obtained by executing the
system on $i$ and observing $p$. Performance prediction approximates the
function $y:C\to\mathbb{R}$ with a regression model $\hat y$ trained on a
sample $S\subset C$ of measured configurations~\cite{guo2013ASE,siegmund2015FSE,kaltenecker2019,julianasurvey}.
Prediction quality is reported on the remaining configurations or on a
designated validation set, typically using the mean relative error
$\MRE = \frac{1}{|C|}\sum_{c\in C}|y_{c}-\hat y_{c}|/y_{c}$~\cite{kaltenecker2019,pereira2020,julianasurvey}.

The dominant axis of comparison between approaches is the trade-off between
prediction accuracy and the effort spent acquiring $S$. The exact
operationalization of this effort is the topic of
\Cref{subsec:cost-models}. The pattern we want to highlight here is that
the strategy chooses what to measure, but the cost term in the literature is
almost always reported as a count, not as the time it actually took.

\subsection{Sampling Strategies}
\label{subsec:sampling-strategies}

We describe the six strategies considered in our experiments. They are the same
six used by Kaltenecker \etal~\cite{kaltenecker2019} and have the practical
advantage of being domain-agnostic, requiring no prior measurements or
knowledge of the system being sampled.

\paragraph*{Random.}
Uniform random sampling draws configurations uniformly at random from $C$. When
$|C|$ is small enough to enumerate, this is straightforward; for larger spaces,
hashing-based samplers such as
\textsc{UniGen}-style approaches~\cite{chakraborty2013scalable} approximate uniformity
under SAT constraints, though Plazar \etal~\cite{plazar2019} have shown that
state-of-the-art uniform samplers face scalability issues. For most systems we
study, $|C|$ is tractable enough that enumeration-based random sampling is the
intended baseline.

\paragraph*{Solver-based.}
Off-the-shelf SAT/SMT solvers (Sat4j~\cite{sat4j}, Z3~\cite{Z3}) generate
configurations by solving the feature model. Internal heuristics make the
exploration order non-uniform but
deterministic~\cite{DBLP:journals/tse/HenardPPKHT14}, so successive calls tend
to produce locally clustered samples.

\paragraph*{Randomized solver-based.}
Henard \etal~\cite{DBLP:journals/tse/HenardPPKHT14} permute variable order and
constraint order on each call, increasing diversity at the cost of any
uniformity guarantee. We use this strategy as in
Kaltenecker \etal~\cite{kaltenecker2019}.

\paragraph*{$t$-wise (coverage-based).}
$t$-wise sampling~\cite{Kuhn2004,Cohen2008,Johansen2012} selects configurations
to cover all combinations of $t$ selected options. We use $t\in\{1,2,3\}$
following Siegmund \etal~\cite{siegmund2012a}, the implementation also used by
Kaltenecker \etal~\cite{kaltenecker2019}. The natural sample size of $t$-wise is
not a free parameter; we record it and include the resulting
$(\text{strategy},n)$ pairs alongside other strategies at common $n$.

\paragraph*{Distance-based.}
Kaltenecker \etal~\cite{kaltenecker2019} propose distance-based sampling, which
selects configurations according to a target distribution over Hamming distance
between configurations. The intent is to approximate the diversity of uniform
random sampling without enumerating $C$.

\paragraph*{Diversified distance-based.}
A variant of distance-based sampling~\cite{kaltenecker2019} that iteratively
adds configurations containing the least frequently selected options, in order
to avoid systematically excluding rare options.

\paragraph*{Diagnostic baselines.}
Beyond these six strategies, we also include simple baselines whose role is
diagnostic rather than competitive: a \emph{cheap-first oracle} that selects
configurations in increasing order of measured cost, and an
\emph{expensive-first oracle} that selects in decreasing order. Neither is
deployable without prior measurements, but together they bound the effect that
cost-awareness could have if a strategy were perfectly cost-informed (cheap
side) or perfectly adversarial (expensive side). They tell us how much
cost-heterogeneity there is to be exploited or stumbled into.

\subsection{Prior Cost Models}
\label{subsec:cost-models}

The closest prior cost model to ours is Sarkar \etal's~\cite{sarkar2015}, which
defines a total cost
\[
\mathit{TotalCost}\;=\;C_{\text{train}} + C_{\text{test}} + C_{\text{model}} + C_{\text{error}},
\]
where $C_{\text{train}}$ and $C_{\text{test}}$ are training and testing
acquisition costs, $C_{\text{model}}$ is model-building cost, and
$C_{\text{error}}$ scales with prediction error. Two simplifications follow.
First, model-building cost is dropped on the grounds that CART training is
inexpensive compared with measurement. Second, $C_{\text{train}}$ and
$C_{\text{test}}$ are collapsed under a $50/50$ train/test split into $2n$,
\emph{treating each measured configuration as having unit cost}. The resulting
formula, $\mathit{TotalCost}(n)=2n+\varepsilon_{n}\cdot|S|\cdot R$, is the
canonical configurable-systems cost model: it is count-based and folds train
and test acquisition together.

Other works adopt similar count-based abstractions, even when they do not
formalize a cost equation. Guo \etal~\cite{guo2013ASE} measure ``feasibility''
in terms of sample size $N,2N,3N,M$. Westermann \etal~\cite{westermann2012}
formulate efficiency as ``required number of measurement points'' to reach a
target error. Siegmund \etal~\cite{siegmund2015FSE} report Pareto trade-offs
between accuracy and ``number of measurements''. Kaltenecker
\etal~\cite{kaltenecker2019} compare strategies at common sample sizes. The
systematic review of Pereira \etal~\cite{julianasurvey} confirms that, when
``cost'' is mentioned in this literature, it almost always means cardinality.

The one configurable-systems paper we are aware of that reports
\emph{empirical} acquisition time is Siegmund \etal~\cite{siegmund2012a}, which
gives both the number of measurements and the time in hours required by each
heuristic. However, that empirical effort is reported \emph{descriptively}; it
is not used to test whether $n$ is a valid proxy, nor to compare strategies at
fixed $n$, nor to study how rankings would change under a time-based
cost. To the best of our knowledge, no prior configurable-systems paper has
done so.

We make two design choices that distinguish our cost formulation from
Sarkar \etal's. First, we restrict the cost term to training: a sampling
strategy controls the configurations it asks to measure, but the validation
set is part of the protocol and would be the same under any strategy. Folding
both into a $2n$ term obscures what the strategy is actually responsible
for. Second, we do not collapse train/test acquisition into a single linear
function of $n$; we compute training cost as the sum of measured
per-configuration times over the configurations the strategy
selects. \Cref{sec:cost_model} formalizes both choices.

\subsection{Context Sensitivity of Performance Prediction}
\label{subsec:context-sensitivity}

A separate strand of work has shown that conclusions about performance
prediction can depend on the input or workload. Pereira
\etal~\cite{pereira2020} showed on x264 that the ranking of sampling
strategies varies markedly across $17$ input videos and two performance
properties. \red{Lesoil et al. (JSS 2023, JSS 2024) and Mühlbauer et al. (ICSE
2023) extend this point to broader collections of configurable systems and to
input-aware models~\cite{}.} Our use of multiple inputs and properties is not
intended to (re)establish that ``context matters''---it is established---but to
ask the narrower question whether the validity of the size-as-cost abstraction
itself is context-dependent: in some inputs $n$ may be a fine proxy and in
others not, and asymmetric cases (predicting output size while measurement
cost is encoding time) are particularly informative.
